% 元の
% \documentclass[lualatex,book,paper=a4paper,jafontsize=11.5pt,jafontscale=1.05]{jlreq}
\documentclass[lualatex,paper=a4paper,jafontsize=11.5pt,jafontscale=1.05]{jlreq}
% 2ページ目や章の前に一ページ余分なものがあるのが気になる場合は openanyを\documentclass[...,book,openany ,..]として追加する

%フォント:
\usepackage{luatexja-fontspec}
\usepackage{hack-fonts}

%表紙用:
\usepackage{hack-cover}
\usepackage{enumitem}
% 画像
\usepackage{graphicx}
\usepackage{float}

% リンク
\usepackage[hidelinks]{hyperref}

% ソースコード
\usepackage{listings}
\lstset{
  frame=single,
  basicstyle=\small\ttfamily,
  tabsize=4,
  breaklines=true
}

% 色・ガントチャート
\usepackage{xcolor}
\usepackage{pgfgantt}
\usepackage{caption}

\begin{document}
\section*{作業ログ}
\noindent
報告者:恩田 隼士\\
報告日: \today \\
プロジェクト名: 単一のフルカラーLEDとArduinoを用いた光無線通信による楽器間同期演奏システム \\
\hrulefill

\section{進捗サマリー（概要）}
\begin{itemize}
 \item 全体進捗率: 95\%
 \item マイルストーン達成状況: 達成
 \item 主要成果:
 \begin{itemize}
    \item 指揮者及び演奏者Arduino間の光無線通信による同期演奏
    \item 複数のArduinoを用いた輪唱演奏の実現
 \end{itemize}
 \item 重大課題（影響度）: なし
\end{itemize}

\section{作業ログ（詳細トラッキング）}
\begin{table}[H]
  \centering
  \caption{作業ログ}
   \scalebox{0.7}{
    \begin{tabular}{lcccccc} \hline
      作業項目 & 開始日 & 予定完了日 & 状態 & 進捗率 & 依存関係・ブロッカー & コメント \\ \hline
      楽器演奏 & 05/20 & 05/23 & 完了 & 100\% & なし & - \\
      演奏開始・停止ボタン（設計） & 05/23 & 05/27 & 完了 & 100\% & 部品入手 & - \\
      楽器音 & 05/23 & 05/27 & 完了 & 100\% & 楽器演奏 & - \\
      楽譜配列 & 05/26 & 05/28 & 完了 & 100\% & 楽器音 & - \\
      ヴァイオリン音 & 05/28 & 06/12 & 完了 & 100\% & 楽器音 & 今週完了 \\
      演奏開始・停止ボタン（実装） & 05/28 & 06/09 & 完了 & 100\% & 設計 & - \\
      結合 & 06/17 & 07/07 & 完了 & 80\% & 進行中 & 複数台のArduinoを用いた輪唱演奏も実現 \\
      開始停止スイッチテスト & 06/13 & 06/19 & 進行中 & 50\% & 実装の完了 & 単体は確認済，結合は他担当待ち \\
      楽器音テスト & 06/13 & 06/19 & 進行中 & 50\% & 実装の完了 & 単体は確認済，結合は他担当待ち \\
      \hline
    \end{tabular}
  }
\end{table}


% ===== 計画前ガントチャート =====
\subsection{担当箇所のガントチャート（計画前）}
\noindent
\begin{figure}[H]
\centering
\resizebox{0.92\linewidth}{!}{%
\begin{ganttchart}[
  hgrid,
  vgrid={*6{draw=none,line width=0pt}, dotted},
  time slot format=isodate,
  time slot unit=day,
  x unit=0.52cm,
  y unit title=0.6cm,
  y unit chart=0.65cm,
  title label font=\tiny,
  bar label font=\small,
  bar label node/.append style={text width=4.8cm, align=right},
  bar height=0.55,
  bar top shift=0.23,
  bar/.append style={fill=violet!45, draw=violet!75},
]{2026-05-20}{2026-06-16}
  \gantttitlecalendar{month=shortname, day} \\
  \ganttbar{楽器演奏}{2026-05-20}{2026-05-23} \\
  \ganttbar{演奏開始・停止ボタン}{2026-05-23}{2026-05-27} \\
  \ganttbar{楽器音}{2026-05-23}{2026-05-27} \\
  \ganttbar{楽譜配列}{2026-05-26}{2026-05-28} \\
  \ganttbar{ヴァイオリン音}{2026-05-28}{2026-05-31}
  \ganttbar[bar label node/.append style={opacity=0}]{}{2026-06-04}{2026-06-09} \\
  \ganttbar{開始停止スイッチテスト}{2026-06-10}{2026-06-16} \\
  \ganttbar{楽器音テスト}{2026-06-10}{2026-06-16}
\end{ganttchart}%
}
\caption{担当箇所ガントチャート（計画前）}
\end{figure}

% ===== 現在のガントチャート =====
% 凡例：■ 灰色 = 完了，■ 紫色 = 未完了 / 進行中
% タスクの状態に応じて fill=gray!60（完了）または fill=violet!50（未完了）を変更
\subsection{担当箇所のガントチャート（現在・進捗状況）}
\noindent
{\small \textcolor{gray!50}{\rule{1em}{0.8em}}\ 完了\quad
       \textcolor{violet!60}{\rule{1em}{0.8em}}\ 未完了／進行中}
\begin{figure}[H]
\centering
\resizebox{0.92\linewidth}{!}{%
\begin{ganttchart}[
  hgrid,
  vgrid={*6{draw=none,line width=0pt}, dotted},
  time slot format=isodate,
  time slot unit=day,
  x unit=0.52cm,
  y unit title=0.6cm,
  y unit chart=0.65cm,
  title label font=\tiny,
  bar label font=\small,
  bar label node/.append style={text width=4.8cm, align=right},
  bar height=0.55,
  bar top shift=0.23,
]{2026-05-20}{2026-06-25}
  \gantttitlecalendar{month=shortname, day} \\
  % ---- 完了タスク (gray!60) ----
  \ganttbar[bar/.append style={fill=gray!60, draw=gray!80}]
    {楽器演奏}{2026-05-20}{2026-05-23} \\
  \ganttbar[bar/.append style={fill=gray!60, draw=gray!80}]
    {演奏開始・停止ボタン（設計）}{2026-05-23}{2026-05-27} \\
  \ganttbar[bar/.append style={fill=gray!60, draw=gray!80}]
    {楽器音}{2026-05-23}{2026-05-27} \\
  \ganttbar[bar/.append style={fill=gray!60, draw=gray!80}]
    {楽譜配列}{2026-05-26}{2026-05-28} \\
  % ---- ヴァイオリン音：今週完了 (gray!60) ----
  \ganttbar[bar/.append style={fill=gray!60, draw=gray!80}]
    {ヴァイオリン音}{2026-05-28}{2026-05-31}
  \ganttbar[bar/.append style={fill=gray!60, draw=gray!80},
    bar label node/.append style={opacity=0}]
    {}{2026-06-04}{2026-06-12} \\
  \ganttbar[bar/.append style={fill=gray!60, draw=gray!80}]
    {演奏開始・停止ボタン（実装）}{2026-05-28}{2026-05-31}
  \ganttbar[bar/.append style={fill=gray!60, draw=gray!80},
    bar label node/.append style={opacity=0}]
    {}{2026-06-04}{2026-06-09} \\
  % ---- 進行中タスク (violet!50)：単体確認済・結合は他担当待ち ----
  \ganttbar[bar/.append style={fill=violet!50, draw=violet!80}]
    {開始停止スイッチテスト}{2026-06-13}{2026-06-19} \\
  \ganttbar[bar/.append style={fill=violet!50, draw=violet!80}]
    {楽器音テスト}{2026-06-13}{2026-06-19} \\
  % ---- 今週追加：実機結合による輪唱演奏 ----
  \ganttbar[bar/.append style={fill=violet!50, draw=violet!80}]
    {輪唱演奏（実機結合）}{2026-06-17}{2026-06-25}
\end{ganttchart}%
}
\caption{担当箇所ガントチャート（現在：灰＝完了，紫＝未完了）}
\end{figure}

ガントチャートにおいて，06/01$\sim$06/03の期間はテスト期間として空白にしている．

% ===== 日ごとの作業時間・内容 =====
\subsection{日ごとの作業時間・内容}
\begin{table}[H]
  \centering
  \caption{日ごとの作業時間・内容}
  \scalebox{0.88}{
    \begin{tabular}{ccp{7.5cm}} \hline
      日付 & 作業時間 & 作業内容・備考 \\ \hline
      2026/06/17 & 4h & 同期演奏システムの結合 \\[4pt]
      2026/06/20 & 6h & Arudinoを複数用いた楽器演奏，歌詞表示と楽器演奏の同期 \\[4pt]
        % ※必要に応じて行を追加
      \hline
    \end{tabular}
  }

\end{table}

% ===== 今週追加：ログの詳細（TPM・TRL） =====
\section{ログの詳細}

\subsection{指揮者及び演奏者Arduino間の光無線通信による同期演奏}
各担当の機能の単体実装が完了したため，実機およびソースコードの結合を行った．
結合対象は指揮者Arduino（\texttt{conductor.ino}），演奏者Arduino（\texttt{player.ino} + \texttt{lyrics.h/cpp} + \texttt{katakana.h}），およびProcessing側音源（\texttt{soundPlayer.pde}）の3層であり，
指揮者LEDの色・点滅で音高・テンポ・音量を伝送し，演奏者がカラーセンサで受光してシリアル経由でProcessingに4Byteパケット（pitch 6bit／volume 7bit／duration 12bit／instID 2bit）を送信する構成である．
結合過程で，演奏者側の色判定が「毎拍判定」となっていた点を共同開発者と認識合わせし，
「一度自色を検出した以降は点滅追従のみで進行する」ラッチ動作に修正した（\texttt{armed}フラグ導入）．
これにより，色を切り替えながら点滅させるだけで複数楽器の同時演奏が成立する状態となった．

\subsection{複数のArduinoを用いた輪唱演奏}
同期演奏システムを拡張し，複数台の演奏者Arduinoによる輪唱（カノン）演奏を実装した．
共通の \texttt{melody[]} を全演奏者で共有したまま，指揮者側の順次点火シーケンスのみで各パートの入りをずらす方式である．
指揮者は OFF→ON 立ち上がりごとに \texttt{pulseIdx} を進め，\texttt{OFFSET\_PULSES\_PER\_PART = 16}（8拍）ごとに点灯色を \texttt{PART\_COLORS[] = \{紫, 赤, 青, 緑\}} の次の色へ切り替える．
全パート点火後は末尾色（緑）で点滅を継続し，最後に入ったパートの色で全演奏者を引き続き同期させる．
これにより，最大4声の輪唱を楽譜・シリアルプロトコルを変更することなく実現した．


\subsection{技術性能指標（TPM）}
TPMはシステムの最終的な性能に関わる技術を管理するための指標である．
本プロジェクトの各指標のうち，自分の担当は「音色の再現性」である．
実楽器音の録音と本システムの出力音を比較する．演奏の正確性（accuracy）は
一意に定義することが難しいため\cite{Larrouy}，本テストでは人間の音色知覚を構成する
以下の4特徴量を定量指標として再現精度を評価する\cite{McAdams}．各楽器・音程ごとに計測し，平均で評価する．

\begin{itemize}
  \item 立ち上がり時間（AT）：RMSエンベロープがピークに到達するまでの時間[ms]．
        目標は全音程平均で60\,ms以内\cite{Siedenburg}．
  \item スペクトル重心（SC）：持続区間のスペクトル重心周波数[Hz]．
        目標は相対差の平均で24\%以内\cite{Wun}．
  \item スペクトル変動度（DSV）：フレーム間のL2相対変化量の平均[\%]．
        目標はズレの平均で16\%以内\cite{McAdams}．
  \item 奇偶倍音比（OER）：pYINアルゴリズムでF0を推定し，奇数次倍音と
        全倍音のエネルギー比を算出．目標は相対差の平均で20\%以内\cite{Caclin}．
\end{itemize}

今週の計測結果（全7音程の平均）を表\ref{tab:tpm}に示す．いずれの特徴量も目標値を満たし，
担当指標である音色の再現性は達成と判定した（詳細は図\ref{fig:violin_result}）．

\begin{table}[H]
  \centering
  \caption{音色の再現性（TPM）計測結果}
  \label{tab:tpm}
  \scalebox{0.9}{
    \begin{tabular}{lcccc} \hline
      特徴量 & 元音平均 & 合成音平均 & 差（目標） & 判定 \\ \hline
      AT [ms]  & 809.4 & 769.5 & 39.6\,ms（$\leq$60\,ms） & 達成 \\
      SC [Hz]  & 2332  & 2492  & 12.2\%（$\leq$24\%） & 達成 \\
      DSV [\%] & 28.6  & 40.5  & 11.67\%（$\leq$16\%） & 達成 \\
      OER [0--1] & 0.746 & 0.746 & 0.097（$\leq$0.20） & 達成 \\
      \hline
    \end{tabular}
  }
\end{table}

\subsection{技術成熟度（TRL）}
TRLは機能ごとの進捗を確認する指標である．定義を表\ref{tab:trl}に示す．

\begin{table}[H]
  \centering
  \caption{TRLの定義}
  \label{tab:trl}
  \scalebox{0.9}{
    \begin{tabular}{cp{10cm}} \hline
      TRL & 定義 \\ \hline
      1 & 概念レベル：アイディアや原理が確立されており，説明できるレベル \\
      2 & 基本動作レベル：単体で動作確認が可能なレベル \\
      3 & 結合動作レベル：機能を結合し，実際の使用環境に近い環境で動作確認が可能なレベル \\
      4 & 実運用レベル：実際の使用環境で安定的な稼働が確認できるレベル \\
      \hline
    \end{tabular}
  }
\end{table}

自分の担当機能の現状を表\ref{tab:trl_status}に示す．同期演奏システムの結合は実機での結合及びテストを進めている段階であるため，TRL3と判定した．
輪唱演奏についても，複数台の演奏者Arduinoと指揮者Arduinoを用いた実機結合での動作確認まで到達しているため，TRL3と判定した．

\begin{table}[H]
  \centering
  \caption{担当機能のTRL現状}
  \label{tab:trl_status}
  \scalebox{0.9}{
    \begin{tabular}{lcp{6cm}} \hline
      担当機能 & 現在のTRL & 根拠 \\ \hline
      同期演奏システムの結合 & TRL3 & 実機での結合及びテストを進めている \\
      輪唱演奏 & TRL3 & 複数台Arduinoでの実機結合動作を確認済 \\
      \hline
    \end{tabular}
  }
\end{table}

\section{課題管理}

\subsection{コード統合}
\begin{itemize}
  \item 原因: 各担当者が個別に実装を進めているため，統合時にコードの調整が必要となる．
  \item 影響度: 中
  \item 優先度: 中
  \item 対応策: 他担当者と実装の仕様・コードを確認し，統合作業を計画的に進める．
  \item 期限: 2026/06/19
  \item 状態: 未着手
\end{itemize}


% ※影響度＝プロジェクトへのインパクト
% ※優先度＝対応の緊急性

\section{AI利用状況と検証}
\subsection{利用内容}
\subsubsection*{A. 演奏者Arduinoの色判定ラッチ化（\texttt{player.ino}）}
\begin{itemize}
    \item 利用目的：合奏成立のため，「一度自色を検出した後は色不問で点滅追従のみ行う」ラッチ動作への修正
    \item 入力（プロンプト要旨）：
    \begin{itemize}
        \item 現状の \texttt{player.ino} は毎拍で \texttt{detectColorID()} を呼ぶため，指揮者LEDの色を切り替えると自パート以外の色で演奏が停止する旨を説明した．
        \item 設計書3.3.5.3および図3.12では「初回自色検出以降は点滅検出のみ」と読み取れる旨を共有し，最小差分でラッチ化する実装を依頼した．
        \item 既存のシリアル送信タイミング・楽譜進行ロジックは変更しないことを条件として指定した．
    \end{itemize}
    \item 出力概要：
    \begin{itemize}
        \item グローバルに \texttt{bool armed = false} を追加し，\texttt{if (!armed)} 内でのみ \texttt{detectColorID()} を呼び，自色一致時に \texttt{armed = true} とする実装が得られた．
        \item 案A（\texttt{noteNum >= MELODY\_LENGTH} で自動停止）のみ採用し，案B（一定時間消灯による再武装）は未実装で保留する判断が得られた．
    \end{itemize}
\end{itemize}

\subsubsection*{B. 指揮者Arduinoの輪唱用順次点火シーケンス（\texttt{conductor.ino}）}
\begin{itemize}
    \item 利用目的：複数演奏者で共通の \texttt{melody[]} を共有したまま，輪唱（カノン）を成立させる点火制御の実装
    \item 入力（プロンプト要旨）：
    \begin{itemize}
        \item 楽譜やシリアルプロトコルは変えず，指揮者LEDの点灯色を一定パルスごとに切り替えるだけで輪唱を成立させたい旨を指定した．
        \item オフセットは8拍，最大4声，末尾に入ったパートの色で点滅継続，という条件を与えた．
        \item \texttt{SYSTEM\_SWITCH} OFF→ON の立ち上がり検出時に \texttt{pulseIdx} をリセットして先頭から再演奏する仕様も指定した．
    \end{itemize}
    \item 出力概要：
    \begin{itemize}
        \item \texttt{NUM\_PARTS = 4}，\texttt{OFFSET\_PULSES\_PER\_PART = 16}（8分音符単位で16パルス），\texttt{PART\_COLORS[] = \{紫, 赤, 青, 緑\}} を定数定義した実装が得られた．
        \item OFF→ON 立ち上がりごとに \texttt{pulseIdx} を進め，\texttt{OFFSET\_PULSES\_PER\_PART} ごとに点灯色を次の色へ切り替え，全パート点火後は末尾色で点滅継続する制御が得られた．
    \end{itemize}
\end{itemize}

\subsubsection*{C. UNO R4 WiFi の \texttt{WiFi.begin()} ブロッキング対策（\texttt{conductor.ino}）}
\begin{itemize}
    \item 利用目的：WiFi接続失敗時の起動待機時間が想定の倍になる事象の解消
    \item 入力（プロンプト要旨）：
    \begin{itemize}
        \item 起動時に「無反応約10秒 → ドット出力約10秒 → スタンドアロン開始」と二段待機する症状を説明した．
        \item 接続失敗時に \texttt{WiFi.begin()} 自体が内部で数秒〜10秒ブロックしてから戻っているのではないかという仮説を提示し，検証と修正案の作成を依頼した．
    \end{itemize}
    \item 出力概要：
    \begin{itemize}
        \item タイマー開始（\texttt{millis()}）を \texttt{WiFi.begin()} の \textbf{前} に置き，\texttt{WiFi.begin()} 自体のブロック時間もタイムアウトに含める修正パターンが得られた．
        \item スタンドアロン起動フォールバック時に \texttt{SYSTEM\_SWITCH = false} を明示セットし，単体モードでは停止状態起動を保証する合わせ技も提案された．
    \end{itemize}
\end{itemize}

\subsection{検証方法}
% ※第三者が再現できるレベルで記述
\begin{itemize}
  \item 検証方法:
  \begin{itemize}
      \item いずれも実機（指揮者Arduino・演奏者Arduino複数台・Processing PC）での結合動作確認による．
  \end{itemize}
  \item 参照資料:
  \begin{itemize}
      \item 本プロジェクト設計書 3.3.5.3 / 図3.12（色判定挙動），同 2.4（状態遷移）．
      \item Arduino UNO R4 WiFi 公式リファレンス（\texttt{WiFiS3} ライブラリ）．
  \end{itemize}
  \item 検証手順：
  \begin{enumerate}
      \item A：指揮者LEDの色を自パート色→他パート色の順に点滅させ，初回自色検出以降は他色でも演奏が継続することを確認する．
      \item B：演奏者Arduinoを複数台接続し，指揮者起動後8拍ごとに新しいパートが入り，全パート点火後は末尾色（緑）で同期して輪唱が継続することを確認する．
      \item C：存在しないSSIDを指定して指揮者Arduinoを起動し，全体待機時間が想定タイムアウト値の範囲内（二段待機にならない）で収まることを確認する．
  \end{enumerate}
\end{itemize}

\subsection{ファクトチェック}
\begin{itemize}
  \item 一致点：
  \begin{itemize}
    \item A：ラッチ動作仕様は，共同開発者（24G1125 三芝）との認識合わせおよび設計書3.3.5.3／図3.12の記述と一致する．
    \item B：順次点火シーケンスは，楽譜・シリアルプロトコル無変更で輪唱を成立させるという当初要件と整合する．
    \item C：\texttt{WiFi.begin()} がリンクアップ失敗時に内部ブロックする挙動は，他のArduino WiFiライブラリでも同様の流派が一般的であり妥当である．
  \end{itemize}
  \item 不一致・疑義：なし
  \item 最終判断：いずれも採用
  \item 根拠：
  \begin{itemize}
    \item A：修正後の \texttt{player.ino} を実機投入し，紫→赤→青→緑と色を切り替えても自パート以外の色で停止しないことを確認した．
    \item B：演奏者Arduinoを複数台接続し，8拍オフセットで順次入る輪唱演奏が成立することを確認した．
    \item C：存在しないSSID指定で起動した際の待機時間が単一タイムアウト分に収まり，二段待機が解消したことを確認した．
  \end{itemize}
\end{itemize}

\section{次回計画（次回までに実施する内容）}
\begin{itemize}
  \item 実施事項：結合テスト（他担当との統合）
  \begin{itemize}
    \item 他担当の実装完了後，ヴァイオリン音・開始停止スイッチを結合し，実使用環境に近い条件で動作テストを実施する（TRL3を目指す）．
  \end{itemize}
  \item 達成基準：
  \begin{itemize}
    \item 結合状態で演奏の開始・停止および各音の発音が正しく動作すること．
  \end{itemize}
  \item 期限：2026/06/19
\end{itemize}

\section{その他}
連絡事項・懸念点：なし

\section{付録}
添付資料を以下のリポジトリURLに示す．添付範囲は全体である．\\
\url{https://github.com/crab2424/LEDplaying_system}



\begin{thebibliography}{9}
\bibitem{processing} Processing Foundation (2026), "Processing Reference". \url{https://processing.org/reference/}. (2026年6月16日 参照)
\bibitem{minim_unpatch} D. Di Fede (2026), "Minim : UGen.unpatch()". \url{http://code.compartmental.net/minim/ugen_method_unpatch.html}. (2026年6月16日 参照)
\bibitem{Vurma} A. Vurma and J. Ross, "Production and perception of musical intervals," \textit{Music Perception}, vol. 23, no. 4, pp. 331--344, 2006.
\bibitem{Larrouy} P. Larrouy-Maestri, "I know it when I hear it," \textit{Music \& Science}, vol. 1, 2018.
\bibitem{McAdams} S. McAdams, S. Winsberg, S. Donnadieu, G. De Soete, and J. Krimphoff, "Perceptual scaling of synthesized musical timbres: Common dimensions, specificities, and latent subject classes," \textit{Psychological Research}, vol. 58, no. 3, pp. 177--192, 1995.
\bibitem{Caclin} A. Caclin, S. McAdams, B. K. Smith, and S. Winsberg, "Acoustic correlates of timbre space dimensions: A confirmatory study using synthetic tones," \textit{The Journal of the Acoustical Society of America}, vol. 118, no. 1, pp. 471--482, 2005.
\bibitem{Siedenburg} K. Siedenburg, "Specifying the perceptual relevance of onset transients for musical instrument identification," \textit{The Journal of the Acoustical Society of America}, vol. 145, no. 2, pp. 1078--1087, 2019.
\bibitem{Wun} S. Wun, A. Horner, and B. Wu, "Effect of spectral centroid manipulation on discrimination and identification of instrument timbres," \textit{Journal of the Audio Engineering Society}, vol. 62, no. 9, pp. 575--583, 2014.
\end{thebibliography}

\end{document}
